% Môn: Toán
% Lớp: 10
% Chương: Chương I. Mệnh đề và tập hợp
% Bài:  Mệnh đề · Mệnh đề tương đương
% Số câu: 0
% App đọc file này trực tiếp — sửa rồi Commit trên GitHub
% Ghi nguồn từng câu: \nguon{SGK} hoặc \nguon{2 SBT VL 12 KNTT} trong


%%%=============EX_1=============%%%

% Mức: NB
% ID: T10C1B1-151-TN
% ID: T10C1B1-151-TN
% Mức: NB
% ID: T10C1B1-151-TN

%%%=============EX_20=============%%%

% Mức: NB
% ID: T10C1B1-152-TN
% ID: T10C1B1-152-TN
% Mức: NB
% ID: T10C1B1-152-TN

%%%=============EX_21=============%%%

% Mức: NB
% ID: T10C1B1-153-TN
% ID: T10C1B1-153-TN
% Mức: NB
% ID: T10C1B1-153-TN

%%%=============EX_22=============%%%

% Mức: NB
% ID: T10C1B1-154-TN
% ID: T10C1B1-154-TN
% Mức: NB
% ID: T10C1B1-154-TN

%%%=============EX_23=============%%%

% Mức: NB
% ID: T10C1B1-155-TN
% ID: T10C1B1-155-TN
% Mức: NB
% ID: T10C1B1-155-TN

%%%=============EX_24=============%%%

% Mức: NB
% ID: T10C1B1-156-TN
% ID: T10C1B1-156-TN
% Mức: NB
% ID: T10C1B1-156-TN
% Mức: NB

%%%=============EX_25=============%%%
\dangbt{Mệnh đề tương đương}
\begin{ex}
	% ID: T10C1B1-156-TN
	% Mức: NB
	Cho tứ giác ${DEGO}$. Xét hai mệnh đề:
	
	$P$: ''Tứ giác ${DEGO}$ là hình thoi''.
	
	$Q$: ''Tứ giác ${DEGO}$ có hai đường chéo vuông góc''.
	
	Mệnh đề nào sau đây là mệnh đề $P\Leftrightarrow Q$?
	\choice
		{ Nếu tứ giác ${DEGO}$ có hai đường chéo vuông góc thì nó không là hình thoi }
		{ Tứ giác ${DEGO}$ có hai đường chéo vuông góc là điều kiện cần để nó là hình thoi }
		{ Tứ giác ${DEGO}$ là hình thoi thì nó có hai đường chéo vuông góc }
		{ \True Tứ giác ${DEGO}$ là hình thoi là điều kiện cần và đủ để tứ giác đó có hai đường chéo vuông góc }
	\loigiai{
		Mệnh đề $P\Leftrightarrow Q$ là mệnh đề tương đương, được phát biểu dưới dạng ''$P$ là điều kiện cần và đủ để có $Q$'' hoặc ''$P$ khi và chỉ khi $Q$''.
		
		Do đó mệnh đề đúng là: ''Tứ giác ${DEGO}$ là hình thoi là điều kiện cần và đủ để tứ giác đó có hai đường chéo vuông góc.''}
\end{ex}
\begin{ex}
% Mức: TH
	Cho hai mệnh đề $P$: "Tam giác $ABC$ đều" và $Q$: "Tam giác $ABC$ cân và có một góc bằng $60^\circ$". Phát biểu mệnh đề tương đương $P \Leftrightarrow Q$ dưới dạng "điều kiện cần và đủ".
	\choice
	{Nếu tam giác $ABC$ đều thì tam giác $ABC$ cân và có một góc bằng $60^\circ$.}
	{Tam giác $ABC$ cân và có một góc bằng $60^\circ$ là điều kiện cần để tam giác $ABC$ đều.}
	{\True Tam giác $ABC$ đều là điều kiện cần và đủ để tam giác $ABC$ cân và có một góc bằng $60^\circ$.}
	{Tam giác $ABC$ đều là điều kiện đủ để tam giác $ABC$ cân và có một góc bằng $60^\circ$.}
	\loigiai{
		Xét định lí dưới dạng tương đương $P \Leftrightarrow Q$, ta có các cách phát biểu sau:
		- "$P$ tương đương $Q$";
		- "$P$ khi và chỉ khi $Q$";
		- "$P$ là điều kiện cần và đủ để có $Q$".
		Ở đây $P$ là "Tam giác $ABC$ đều" và $Q$ là "Tam giác $ABC$ cân và có một góc bằng $60^\circ$". Do đó, mệnh đề tương đương phát biểu dưới dạng điều kiện cần và đủ là: "Tam giác $ABC$ đều là điều kiện cần và đủ để tam giác $ABC$ cân và có một góc bằng $60^\circ$".
	}
	\nguon{SGK KNTT Toán Lớp 10 Bài 1 trang 9 Luyện tập 4}
\end{ex}

\begin{ex}
% Mức: TH
	Cho tam giác $ABC$ có trung tuyến $AM$. Xét hai mệnh đề $P$: "Tam giác $ABC$ vuông tại $A$" và $Q$: "Độ dài đường trung tuyến $AM$ bằng nửa độ dài cạnh $BC$". Phát biểu nào sau đây là mệnh đề tương đương đúng?
	\choice
	{Nếu tam giác $ABC$ vuông tại $A$ thì độ dài đường trung tuyến $AM$ bằng nửa cạnh $BC$.}
	{Tam giác $ABC$ vuông tại $A$ là điều kiện đủ để độ dài đường trung tuyến $AM$ bằng nửa cạnh $BC$.}
	{\True Tam giác $ABC$ vuông tại $A$ khi và chỉ khi độ dài đường trung tuyến $AM$ bằng nửa cạnh $BC$.}
	{Tam giác $ABC$ vuông tại $A$ là điều kiện cần để độ dài đường trung tuyến $AM$ bằng nửa cạnh $BC$.}
	\loigiai{
		Cả hai mệnh đề kéo theo $P \Rightarrow Q$ và $Q \Rightarrow P$ đều đúng theo các định lí hình học lớp 8. Do đó mệnh đề tương đương $P \Leftrightarrow Q$ là mệnh đề đúng, được phát biểu dưới dạng tương đương "khi và chỉ khi": "Tam giác $ABC$ vuông tại $A$ khi và chỉ khi độ dài đường trung tuyến $AM$ bằng nửa cạnh $BC$".
	}
	\nguon{Toán 10 Học kì 1 GV.pdf Bài 1 trang 10 Bài 5}
\end{ex}

\begin{ex}
% Mức: TH
	Cho hai số thực $x, y$. Mệnh đề tương đương nào sau đây là đúng?
	\choice
	{$x > 0 \Leftrightarrow x^2 > 0$}
	{$x > y \Leftrightarrow x^2 > y^2$}
	{\True $x^2 + y^2 = 0 \Leftrightarrow x = 0 \text{ và } y = 0$}
	{$x^2 = y^2 \Leftrightarrow x = y$}
	\loigiai{
		- Xét phương án C: Với mọi số thực $x, y$ ta luôn có $x^2 \ge 0$ và $y^2 \ge 0$. Do đó, $x^2 + y^2 = 0$ khi và chỉ khi đồng thời $x^2 = 0$ và $y^2 = 0 \Leftrightarrow x = 0$ và $y = 0$. Đẳng thức tương đương này hoàn toàn đúng.
		- Các phương án khác đều sai vì thiếu các trường hợp số âm (chẳng hạn với $x = -2$ thì $x^2 = 4 > 0$ nhưng $x$ không lớn hơn $0$).
	}
	\nguon{SBT KNTT Toán Lớp 10 Bài 1 Câu 1.6a}
\end{ex}

\begin{ex}
% Mức: TH
	Cho số tự nhiên $n$. Mệnh đề tương đương nào sau đây là đúng?
	\choice
	{$n$ chia hết cho $5 \Leftrightarrow n$ có chữ số tận cùng bằng $5$}
	{\True $n$ chia hết cho $6 \Leftrightarrow n$ chia hết cho cả $2$ và $3$}
	{$n$ chia hết cho $3 \Leftrightarrow n$ chia hết cho $9$}
	{$n$ chia hết cho $2 \Leftrightarrow n$ có chữ số tận cùng bằng $4$}
	\loigiai{
		- Vì $2$ và $3$ là hai số nguyên tố cùng nhau có tích là $6$, nên một số tự nhiên $n$ chia hết cho $6$ khi và chỉ khi số đó chia hết cho cả $2$ và $3$. Đây là mệnh đề tương đương đúng.
		- Phương án A sai vì số tận cùng là $0$ vẫn chia hết cho $5$.
		- Phương án C sai vì $15$ chia hết cho $3$ nhưng không chia hết cho $9$.
		- Phương án D sai vì số tận cùng là $0, 2, 6, 8$ vẫn chia hết cho $2$.
	}
	\nguon{Toán 10 Học kì 1 GV.pdf Bài 1 trang 9}
\end{ex}

\begin{ex}
% Mức: TH
	Cho hai mệnh đề $P$: "Tứ giác $ABCD$ là hình vuông" và $Q$: "Tứ giác $ABCD$ là hình chữ nhật có hai đường chéo vuông góc với nhau". Nhận xét nào sau đây là đúng về mối quan hệ giữa $P$ và $Q$?
	\choice
	{Mệnh đề $P \Rightarrow Q$ đúng nhưng mệnh đề đảo $Q \Rightarrow P$ sai.}
	{Mệnh đề $Q \Rightarrow P$ đúng nhưng mệnh đề đảo $P \Rightarrow Q$ sai.}
	{\True Cả hai mệnh đề $P \Rightarrow Q$ và $Q \Rightarrow P$ đều đúng, do đó mệnh đề tương đương $P \Leftrightarrow Q$ đúng.}
	{Cả hai mệnh đề $P \Rightarrow Q$ và $Q \Rightarrow P$ đều sai.}
	\loigiai{
		- Theo định nghĩa và tính chất hình học, hình vuông là hình chữ nhật có hai đường chéo vuông góc với nhau, và ngược lại hình chữ nhật có hai đường chéo vuông góc cũng chính là hình vuông. Do đó, cả hai mệnh đề kéo theo đều đúng, dẫn tới mệnh đề tương đương $P \Leftrightarrow Q$ là mệnh đề đúng.
	}
	\nguon{SGK KNTT Toán Lớp 10 Bài 1 trang 9 Ví dụ 5}
\end{ex}

\begin{ex}
% Mức: TH
	Xét số tự nhiên $n$. Mệnh đề tương đương nào sau đây là một mệnh đề sai?
	\choice
	{$n$ là số chẵn $\Leftrightarrow n$ chia hết cho $2$}
	{$n$ chia hết cho $9 \Leftrightarrow$ tổng các chữ số của $n$ chia hết cho $9$}
	{\True $n$ chia hết cho $5 \Leftrightarrow n$ có chữ số tận cùng bằng $5$}
	{$n$ chia hết cho cả $3$ và $5 \Leftrightarrow n$ chia hết cho $15$}
	\loigiai{
		- Mệnh đề ở phương án C sai vì chiều nghịch đúng (một số có tận cùng bằng 5 thì chia hết cho 5) nhưng chiều thuận sai (số tự nhiên chia hết cho 5 có thể có chữ số tận cùng bằng 0, ví dụ số 10). Do đó mệnh đề tương đương này là một mệnh đề sai.
		- Các mệnh đề tương đương ở các phương án A, B, D đều là các định lí chia hết đúng đã biết.
	}
	\nguon{SGK KNTT Toán Lớp 10 Bài 1 trang 9 HĐ6}
\end{ex}

\begin{ex}
% Mức: TH
	Trong các mệnh đề tương đương dưới đây, mệnh đề nào là một mệnh đề đúng?
	\choice
	{Tứ giác $ABCD$ là hình chữ nhật $\Leftrightarrow$ tứ giác $ABCD$ có hai đường chéo bằng nhau.}
	{Tứ giác $ABCD$ là hình thoi $\Leftrightarrow$ tứ giác $ABCD$ có hai đường chéo vuông góc với nhau.}
	{\True Tam giác $ABC$ cân tại $A \Leftrightarrow$ tam giác $ABC$ có hai đường trung tuyến tương ứng bằng nhau.}
	{Tứ giác $ABCD$ là hình vuông $\Leftrightarrow$ tứ giác $ABCD$ có bốn cạnh bằng nhau.}
	\loigiai{
		- Phương án A sai vì hình thang cân cũng có hai đường chéo bằng nhau nhưng không phải là hình chữ nhật.
		- Phương án B sai vì tứ giác có hai đường chéo vuông góc chưa chắc là hình thoi (phải cắt nhau tại trung điểm mỗi đường).
		- Phương án D sai vì hình thoi có bốn cạnh bằng nhau nhưng không phải hình vuông.
		- Phương án C đúng theo định lí về dấu hiệu nhận biết tam giác cân (tam giác có hai đường trung tuyến bằng nhau là tam giác cân).
	}
	\nguon{Toán 10 Học kì 1 GV.pdf Bài 1 trang 9}
\end{ex}

\begin{ex}
% Mức: TH
	Cho hai số thực $a, b$. Mệnh đề tương đương nào sau đây biểu thị đúng điều kiện để tích $ab$ bằng $0$?
	\choice
	{$ab = 0 \Leftrightarrow a = 0 \text{ và } b = 0$}
	{\True $ab = 0 \Leftrightarrow a = 0 \text{ hoặc } b = 0$}
	{$ab = 0 \Leftrightarrow a = 0$}
	{$ab = 0 \Leftrightarrow b = 0$}
	\loigiai{
		Tích của hai số thực bằng $0$ khi và chỉ khi ít nhất một trong hai số đó bằng $0$. Do đó, mệnh đề tương đương đúng phải là: $ab = 0 \Leftrightarrow a = 0$ hoặc $b = 0$.
	}
	\nguon{Toán 10 Học kì 1 GV.pdf Bài 1 trang 12}
\end{ex}

\begin{ex}
% Mức: TH
	Cho tam giác $ABC$. Xét hai mệnh đề $P$: "Tam giác $ABC$ đều" và $Q$: "Tam giác $ABC$ cân và có một góc bằng $60^\circ$". Xét tính đúng sai của các khẳng định sau:
	\choiceTF
	{\True Mệnh đề kéo theo $P \Rightarrow Q$ là một mệnh đề đúng.}
	{\True Mệnh đề kéo theo $Q \Rightarrow P$ là một mệnh đề đúng.}
	{\True Mệnh đề tương đương $P \Leftrightarrow Q$ là một mệnh đề đúng.}
	{Cụm từ "$P$ là điều kiện cần và đủ để có $Q$" dùng để phát biểu cho mệnh đề đảo $Q \Rightarrow P$.}
	\loigiai{
		- Ý a) Đúng vì tam giác đều thì hiển nhiên cân và có góc $60^\circ$.
		- Ý b) Đúng vì tam giác cân có một góc bằng $60^\circ$ là tam giác đều.
		- Ý c) Đúng vì cả hai mệnh đề kéo theo đều đúng, dẫn đến mệnh đề tương đương $P \Leftrightarrow Q$ đúng.
		- Ý d) Sai vì cụm từ "điều kiện cần và đủ" dùng để phát biểu cho mệnh đề tương đương $P \Leftrightarrow Q$ chứ không phải cho mệnh đề đảo.
	}
	\nguon{SGK KNTT Toán Lớp 10 Bài 1 trang 9 Luyện tập 4}
\end{ex}

\begin{ex}
% Mức: TH
	Cho hai số thực $x, y$. Xét tính đúng sai của các mệnh đề tương đương dưới đây:
	\choiceTF
	{$x^2 > 0 \Leftrightarrow x > 0$.}
	{\True $x^2 + y^2 = 0 \Leftrightarrow x = 0$ và $y = 0$.}
	{$xy > 0 \Leftrightarrow x > 0$ và $y > 0$.}
	{$x = y \Leftrightarrow x^2 = y^2$.}
	\loigiai{
		- Ý a) Sai vì với $x = -2$, ta có $(-2)^2 = 4 > 0$ nhưng không có $x > 0$.
		- Ý b) Đúng vì $x^2 \ge 0$ và $y^2 \ge 0$ với mọi $x, y$, do đó tổng bằng 0 khi và chỉ khi từng số bằng 0.
		- Ý c) Sai vì $xy > 0$ còn có trường hợp cả hai số cùng âm (ví dụ $x = -1, y = -2$).
		- Ý d) Sai vì với $x = -2, y = 2$ thì $x^2 = y^2 = 4$ nhưng $x \neq y$.
	}
	\nguon{SBT KNTT Toán Lớp 10 Bài 1 Câu 1.6}
\end{ex}

\begin{ex}
% Mức: TH
	Cho hai mệnh đề $P$: "Tam giác $ABC$ đều" và $Q$: "Tam giác $ABC$ có ba góc bằng nhau". Mệnh đề tương đương $P \Leftrightarrow Q$ có giá trị logic là Đúng hay Sai? (Điền "Đúng" hoặc "Sai").
	\shortans{Đúng}
	\loigiai{
		- Chiều thuận: "Nếu tam giác $ABC$ đều thì tam giác $ABC$ có ba góc bằng nhau" (đều bằng $60^\circ$) là khẳng định đúng.
		- Chiều đảo: "Nếu tam giác $ABC$ có ba góc bằng nhau thì tam giác $ABC$ đều" là khẳng định đúng.
		Vì cả hai chiều kéo theo đều đúng nên mệnh đề tương đương $P \Leftrightarrow Q$ là mệnh đề đúng.
	}
	\nguon{SGK KNTT Toán Lớp 10 Bài 1 trang 9}
\end{ex}

\begin{ex}
% Mức: TH
	Cho hai số thực $x, y$. Cho mệnh đề tương đương $P \Leftrightarrow Q$ với $P: "x^2 + y^2 = 0"$ và $Q: "x = 0 \text{ và } y = 0"$. Xét tính đúng sai của mệnh đề tương đương này (Điền "Đúng" hoặc "Sai").
	\shortans{Đúng}
	\loigiai{
		- Vì $x^2 \ge 0$ và $y^2 \ge 0$ với mọi số thực $x, y$, nên tổng $x^2 + y^2$ luôn lớn hơn hoặc bằng $0$.
		- Đẳng thức $x^2 + y^2 = 0$ chỉ xảy ra khi và chỉ khi đồng thời $x^2 = 0$ và $y^2 = 0 \Leftrightarrow x = 0$ và $y = 0$.
		Do đó mệnh đề tương đương $P \Leftrightarrow Q$ là mệnh đề đúng.
	}
	\nguon{SBT KNTT Toán Lớp 10 Bài 1 Câu 1.6a}
\end{ex}

\begin{ex}
% Mức: TH
	Cho số tự nhiên $n$. Xét mệnh đề tương đương: "Số tự nhiên $n$ chia hết cho $9$ khi và chỉ khi tổng các chữ số của $n$ chia hết cho $9$". Mệnh đề này là đúng hay sai? (Điền "Đúng" hoặc "Sai").
	\shortans{Đúng}
	\loigiai{
		Theo dấu hiệu chia hết cho 9 đã học ở cấp hai, một số tự nhiên chia hết cho 9 khi và chỉ khi tổng các chữ số của nó chia hết cho 9. Do đó đây là một mệnh đề tương đương đúng.
	}
	\nguon{Toán 10 Học kì 1 GV.pdf Bài 1 trang 11}
\end{ex}

\begin{ex}
% Mức: TH
	Cho tam giác $ABC$ với đường trung tuyến $AM$. Xét các mệnh đề:
	$P$: "Tam giác $ABC$ vuông tại $A$";
	$Q$: "Độ dài đường trung tuyến $AM$ bằng nửa độ dài cạnh $BC$".
	Hãy phát biểu mệnh đề tương đương $P \Leftrightarrow Q$ và xét tính đúng sai của nó.
	\loigiai{
		- Phát biểu mệnh đề tương đương $P \Leftrightarrow Q$: "Tam giác $ABC$ vuông tại $A$ khi và chỉ khi độ dài đường trung tuyến $AM$ bằng nửa độ dài cạnh $BC$".
		- Xét tính đúng sai của mệnh đề tương đương này:
		+ Chiều thuận $P \Rightarrow Q$: "Nếu tam giác $ABC$ vuông tại $A$ thì độ dài trung tuyến $AM$ ứng với cạnh huyền bằng nửa cạnh huyền $BC$". Đây là tính chất hình học cơ bản của tam giác vuông, do đó $P \Rightarrow Q$ là mệnh đề đúng.
		+ Chiều đảo $Q \Rightarrow P$: "Nếu tam giác $ABC$ có độ dài trung tuyến $AM$ bằng nửa cạnh tương ứng $BC$ thì tam giác $ABC$ vuông tại $A$". Đây cũng là một dấu hiệu nhận biết tam giác vuông đã được chứng minh, do đó $Q \Rightarrow P$ là mệnh đề đúng.
		Vì cả hai chiều kéo theo đều đúng nên mệnh đề tương đương $P \Leftrightarrow Q$ là mệnh đề đúng.
	}
	\nguon{Toán 10 Học kì 1 GV.pdf Bài 1 Bài 5}
\end{ex}

\begin{ex}
% Mức: TH
	Cho hai số thực $x, y$. Xét các mệnh đề sau:
	$P: "x^2 + y^2 = 0"$;
	$Q: "x = 0 \text{ và } y = 0"$.
	Hãy phát biểu mệnh đề $P \Leftrightarrow Q$ và chứng minh mệnh đề tương đương này là đúng.
	\loigiai{
		- Phát biểu mệnh đề tương đương $P \Leftrightarrow Q$: "$x^2 + y^2 = 0$ khi và chỉ khi $x = 0$ và $y = 0$".
		- Chứng minh tính đúng đắn của mệnh đề tương đương:
		+ Chứng minh chiều thuận $P \Rightarrow Q$: Giả sử ta có $x^2 + y^2 = 0$. Vì bình phương của một số thực luôn không âm nên $x^2 \ge 0$ và $y^2 \ge 0$ với mọi $x, y$. Suy ra tổng $x^2 + y^2 \ge 0$. Đẳng thức xảy ra khi và chỉ khi đồng thời $x^2 = 0$ và $y^2 = 0 \Leftrightarrow x = 0$ và $y = 0$. Vậy chiều thuận đúng.
		+ Chứng minh chiều nghịch $Q \Rightarrow P$: Nếu $x = 0$ và $y = 0$ thì thế trực tiếp vào vế trái ta có $x^2 + y^2 = 0^2 + 0^2 = 0$. Vậy chiều nghịch đúng.
		Vì cả hai chiều đều đúng nên mệnh đề tương đương $P \Leftrightarrow Q$ là đúng.
	}
	\nguon{SBT KNTT Toán Lớp 10 Bài 1 Câu 1.6a}
\end{ex}

\begin{ex}
% Mức: TH
	Cho hai mệnh đề:
	$P$: "Tứ giác $ABCD$ là hình thoi";
	$Q$: "Tứ giác $ABCD$ có hai đường chéo vuông góc với nhau".
	Hãy phát biểu mệnh đề tương đương $P \Leftrightarrow Q$ dưới dạng điều kiện cần và đủ và xét tính đúng sai của phát biểu đó.
	\loigiai{
		- Phát biểu mệnh đề tương đương $P \Leftrightarrow Q$ dưới dạng điều kiện cần và đủ: "Tứ giác $ABCD$ là hình thoi là điều kiện cần và đủ để tứ giác $ABCD$ có hai đường chéo vuông góc với nhau".
		- Xét tính đúng sai của phát biểu này:
		+ Xét chiều thuận $P \Rightarrow Q$: "Nếu tứ giác $ABCD$ là hình thoi thì hai đường chéo vuông góc với nhau". Đây là một tính chất đúng của hình thoi, nên $P \Rightarrow Q$ là mệnh đề đúng.
		+ Xét chiều nghịch $Q \Rightarrow P$: "Nếu tứ giác $ABCD$ có hai đường chéo vuông góc với nhau thì tứ giác $ABCD$ là hình thoi". Đây là khẳng định sai, ví dụ một tứ giác có hai đường chéo vuông góc nhưng chúng không cắt nhau tại trung điểm của mỗi đường thì tứ giác đó không phải là hình thoi.
		Vì mệnh đề kéo theo đảo $Q \Rightarrow P$ là mệnh đề sai, nên mệnh đề tương đương $P \Leftrightarrow Q$ là một mệnh đề sai.
	}
	\nguon{SGK KNTT Toán Lớp 10 Bài 1 trang 6 Ví dụ 3}
\end{ex}

\begin{ex}
% Mức: TH
	Hãy phát biểu điều kiện cần và đủ để số tự nhiên $n$ chia hết cho $2$, sử dụng thuật ngữ "số chẵn" và giải thích tại sao phát biểu đó đúng.
	\loigiai{
		- Phát biểu mệnh đề tương đương: "Số tự nhiên $n$ là một số chẵn là điều kiện cần và đủ để số tự nhiên $n$ chia hết cho $2$".
		- Giải thích tính đúng đắn:
		+ Chiều thuận: "Nếu số tự nhiên $n$ là số chẵn thì $n$ chia hết cho $2$". Định nghĩa số chẵn là số chia hết cho $2$, do đó chiều thuận luôn đúng.
		+ Chiều nghịch: "Nếu số tự nhiên $n$ chia hết cho $2$ thì $n$ là số chẵn". Theo định lí toán học, số tự nhiên chia hết cho $2$ có chữ số tận cùng là một trong các số $0, 2, 4, 6, 8$, do đó nó là số chẵn. Chiều nghịch đúng.
		Vì hai chiều thuận nghịch đều đúng nên mệnh đề tương đương phát biểu dưới dạng điều kiện cần và đủ trên là hoàn toàn chính xác.
	}
	\nguon{SGK KNTT Toán Lớp 10 Bài 1 trang 9 Luyện tập 4}
\end{ex}